\documentclass[a4,11pt]{aleph-notas}

% -- Paquetes adicionales
\usepackage{enumitem}
\usepackage{aleph-comandos}
\usepackage{systeme}

% -- Datos  
\institucion{Facultad de Ciencias Exactas, Naturales y Ambientales}
\carrera{Ciencia de datos}
\asignatura{Álgebra lineal}
\tema{Material extra no. 6: Aplicaciones Lineales}
\autor{Andrés Merino}
\fecha{Periodo 2026-1}

\logouno[0.16\textwidth]{Logos/logoPUCE_04_ac}
\definecolor{colortext}{HTML}{1957d1}
\definecolor{colordef}{HTML}{1957d1}
\fuente{montserrat}

% -- Comandos adicionales
\setlist[enumerate]{label=\roman*.}

\begin{document}

\encabezado

\begin{ejer}
	Sea $\func{T}{\R^4}{\Mat{2}{2}}$ la aplicación lineal definida por
	\[
		T(a,b,c,d)=
		\begin{pmatrix}
			a+b & b+c\\
			c+d & d+a
		\end{pmatrix}.
	\]
	Calcule la nulidad de $T$.
\end{ejer}

\begin{proof}[Solución]
	Para calcular la nulidad de $T$, determinamos primero el kernel. Tomemos $(a,b,c,d)\in\R^4$ tal que
	\[
		T(a,b,c,d)=
		\begin{pmatrix}
			0 & 0\\
			0 & 0
		\end{pmatrix},
	\]
    es decir,
    \[
        \begin{pmatrix}
			a+b & b+c\\
			c+d & d+a
		\end{pmatrix}=
        \begin{pmatrix}
            0 & 0\\
            0 & 0
        \end{pmatrix}.
    \]
	lo cual equivale al sistema de ecuaciones
	\[
		\systeme{
		a+b = 0{,},
		b+c = 0{,},
		c+d = 0{,},
		d+a = 0{.}
		}
	\]
	La matriz aumentada asociada al sistema es
	\[
		\left(
		\begin{array}{cccc|c}
			1 & 1 & 0 & 0 & 0\\
			0 & 1 & 1 & 0 & 0\\
			0 & 0 & 1 & 1 & 0\\
			1 & 0 & 0 & 1 & 0
		\end{array}
		\right)
		\sim
		\left(
		\begin{array}{cccc|c}
			1 & 0 & 0 & 1 & 0\\
			0 & 1 & 0 & -1 & 0\\
			0 & 0 & 1 & 1 & 0\\
			0 & 0 & 0 & 0 & 0
		\end{array}
		\right).
	\]
	De aquí obtenemos que
	\[
		a=-d,\qquad b=d,\qquad c=-d.
	\]
	Si ponemos $d=t$ con $t\in\R$, entonces la solución general del sistema es
	\[
		(a,b,c,d)=(-t,t,-t,t).
	\]
	Por lo tanto,
	\[
		\ker(T)=\{(-t,t,-t,t): t\in\R\}.
	\]

    Ahora, busquemos una base para $\ker(T)$, tomemos un elemento y tenemos que
    \[
        (-t,t,-t,t) = t (-1,1,-1,1)
    \]
	Así, una posible base de $\ker(T)$ es
	\[
		B=\{(-1,1,-1,1)\}.
	\]
	Como esta base tiene un solo vector, el linealmente independiente, se concluye que
	\[
		\dim(\ker(T))=1.
	\]
	Por definición, la nulidad de $T$ es
	\[
		\operatorname{nul}(T)=\dim(\ker(T))=1.
	\]
	En consecuencia, la nulidad de $T$ es $1$.\qedhere
\end{proof}

\begin{ejer}
	Sea $\func{T}{\R^4}{\Mat{2}{2}}$ la aplicación lineal definida por
	\[
		T(a,b,c,d)=
		\begin{pmatrix}
			a+b & b+c\\
			c+d & d+a
		\end{pmatrix}.
	\]
	Calcule el rango de $T$.
\end{ejer}

\begin{proof}[Solución]
	Para calcular el rango de $T$, determinamos primero la imagen. Tomemos $\begin{pmatrix}x & y\\
z & w
\end{pmatrix}\in\Mat{2}{2}$ tal que
	\[
		T(a,b,c,d)=
		\begin{pmatrix}
			x & y\\
            z & w
		\end{pmatrix},
	\]
    es decir,
    \[
        \begin{pmatrix}
			a+b & b+c\\
			c+d & d+a
		\end{pmatrix}=
        \begin{pmatrix}
            x & y\\
            z & w
        \end{pmatrix}.
    \]
	lo cual equivale al sistema de ecuaciones
	\[
		\systeme{
		a+b = x{,},
		b+c = y{,},
		c+d = z{,},
		d+a = w{.}
		}
	\]
	La matriz aumentada asociada al sistema es
	\[
		\left(
		\begin{array}{cccc|c}
			1 & 1 & 0 & 0 & x\\
			0 & 1 & 1 & 0 & y\\
			0 & 0 & 1 & 1 & z\\
			1 & 0 & 0 & 1 & w
		\end{array}
		\right)
		\sim
		\left(
		\begin{array}{cccc|c}
			1 & 1 & 0 & 0 & x\\
			0 & 1 & 1 & 0 & y\\
			0 & 0 & 1 & 1 & z\\
			0 & 0 & 0 & 0 & w-x+y-z
		\end{array}
		\right).
	\]
	De aquí, para que el sistema sea consistente, se debe cumplir que $w-x+y-z=0$, es decir, $w=x-y+z$. Por lo tanto, la imagen de $T$ es
    \[
        \img(T)=\left\{\begin{pmatrix}x & y\\
            z & w
            \end{pmatrix}\in\Mat{2}{2}: w=x-y+z\right\}.
    \]
    Podemos escribir cada elemento de la imagen como
    \[
        \begin{pmatrix}x & y\\
            z & w
            \end{pmatrix}=
        \begin{pmatrix}x & y \\
            z & x-y+z
            \end{pmatrix}=
        x\begin{pmatrix}1 & 0\\
            0 & 1
            \end{pmatrix}+
        y\begin{pmatrix}0 & 1\\
            0 & -1
            \end{pmatrix}+
        z\begin{pmatrix}0 & 0\\
            1 & 1
            \end{pmatrix}.
    \]
    Así, una posible base de $\img(T)$ es
    \[
        B=\left\{\begin{pmatrix}1 & 0\\
            0 & 1
            \end{pmatrix},\begin{pmatrix}0 & 1\\
            0 & -1
            \end{pmatrix},\begin{pmatrix}0 & 0\\
            1 & 1
            \end{pmatrix}\right\}.
    \]
    Este conjunto es linealmente independiente (se debería verificar), por lo que se concluye que
    \[
        \dim(\img(T))=3.
    \]
    Por definición, el rango de $T$ es
    \[
        \operatorname{rango}(T)=\dim(\img(T))=3.
    \]
    En consecuencia, el rango de $T$ es $3$.\qedhere
\end{proof}

\end{document}